\section{The invariant norm of the definite difference}

% BEGIN A75B1
\paragraph{A75B1. Fourier support from vanishing in rational directions.}
Put $h=N+1>0$, $s=|\Psi|$, $Y_\tau=\operatorname{Im}Z_\tau$ and
$D(Y)=\prod_\tau\det Y_\tau$. Suppose $f(Z)$ is holomorphic on the
product of the Hermitian half spaces, periodic under the translations
by $R$ times the integral Hermitian matrices for a fixed positive
integer $R$, and bounded by $B$ for all $Y_\tau\ge (\delta/4)I$,
where $\delta>0$. Suppose for each nonzero $v\in E^h$ and every
period $Z$ that $f(Z+it(v_\tau v_\tau^*)_\tau)\to0$ as
$t\to+\infty$. Then for each $b\ge0$,
\[
  \sup_{Y_\tau\ge\delta I}|f(X+iY)|D(Y)^b\le C B,
\]
where $C$ depends only on the translation lattice, $h,\Psi,\delta,b$.

Indeed, the integral Hermitian matrices form a full lattice in
$\prod_\tau\operatorname{Herm}_h(\mathbb C)$ by the Minkowski
embeddings of the integers of $E$ off the diagonal and of $E^+$ on
the diagonal. Its rational span is the Hermitian matrices over $E$.
The pairing $[T,X]=\sum_\tau\operatorname{tr}(T_\tau X_\tau)$
is rational on that span (it is the real-field trace of the matrix
trace) and nondegenerate since $[T,T]>0$ for $T\ne0$. Thus every
frequency in the dual period lattice is a rational Hermitian $T$.
A single positive integer $c$ clears the entries of a basis of this
dual, hence all its entries.

The Fourier coefficient of $x\mapsto f(x+iY)$ at $T$ equals
$a_T e^{-2\pi[T,Y]}$ with $a_T$ independent of $Y$: differentiate
the integral over the period torus along a Hermitian direction and
integrate by parts using the Cauchy--Riemann equations. The Fourier
series equals $f$ at each such slice, with absolute convergence by
repeated integration by parts on the compact torus. In particular
$|a_T|\le B e^{2\pi[T,Y]}$ for $Y_\tau\ge(\delta/4)I$.
If one $T_\tau$ had a negative direction, increasing $Y$ along the
corresponding positive rank-one matrix at that place would force
$a_T=0$. For a singular $T$ choose a nonzero $v\in E^h$ in its
kernel (singularity at any place means $\det T=0$ in $E$). Along
$Y(t)=(\delta/4)I+t(v_\tau v_\tau^*)_\tau$, the exponential
magnitude in the coefficient integral is constant, while $f$
tends pointwise to zero. The constant bound on the finite torus
allows dominated convergence, hence again $a_T=0$.

For every surviving $T$, all $T_\tau$ are positive definite and
$\prod_\tau\det T_\tau\ge c^{-hs}$, because the norm of the totally
positive integer $\det(cT)$ in $E^+$ is at least one. Write
$P=\sum_\tau\operatorname{tr} T_\tau$ and $Q=[T,Y]$.
The arithmetic--geometric mean inequality on the $hs$ positive
eigenvalues of $Y_\tau^{1/2}T_\tau Y_\tau^{1/2}$ gives
$Q\ge (hs/c)D(Y)^{1/(hs)}$; also $Q\ge\delta P$.
Taking the coefficient bound at $(\delta/4)I$ bounds each absolute
Fourier term on the higher cone by
\[
 B\exp[-\pi(hs/c)D(Y)^{1/(hs)}]\exp[-\pi\delta P/2].
\]
The latter exponential is summable over the positive frequencies:
$P$ dominates the Euclidean norm of their tuple of matrices, and
a lattice has polynomially many points in a ball. The first
exponential times any fixed power $D(Y)^b$ is uniformly bounded.
This proves the assertion.
% END A75B1

% BEGIN A75B2
\paragraph{A75B2. Bounded-denominator reduction into one high cone.}
There exist $c\in\mathbb Z_{>0}$ and $\delta>0$ depending only on
$E,h$ such that every period $Z$ admits a rational unitary matrix
$g$ with $cg$ integral and $\operatorname{Im}(gZ)_\tau\ge\delta I$
for all $\tau$.

Apply Borel's rational reduction theorem to the connected reductive
$\mathbb Q$-group $G=\operatorname{Res}_{E^+/\mathbb Q}U(J)$,
where $J$ is the standard split skew-Hermitian matrix on $E^{2h}$.
Its real points are a product of $U(h,h)$ and its arithmetic subgroup
$\Gamma$ is the stabilizer of the integral lattice $\mathcal O_E^{2h}$
in the defining rational representation. Choose the rational minimal
parabolic $P$ stabilizing its standard full isotropic flag, and the
maximal compact $K$ stabilizing the tuple $iI$. Its compatible
maximal rational split torus has positive part
$\operatorname{diag}(a_1,\ldots,a_h,a_1^{-1},\ldots,a_h^{-1})$,
with $a_j>0$ common to the CM places. The upper triangular changes of
the isotropic basis give simple roots $a_i/a_{i+1}$, and the final
Hermitian translation gives $a_h^2$. Borel's covering gives
\[
 G(\mathbb R)=\Gamma\,\mathcal C\,\Omega A_t K,
\]
with $\mathcal C\subset G(\mathbb Q)$ finite, $\Omega\subset P(\mathbb R)$
compact, and all the displayed simple roots at least $t>0$ on $A_t$.
We use the same rational covering theorem as in the reader's proof
of row integrability. See A.~Borel, \emph{Introduction aux groupes
arithm\'etiques} (1969), Th\'eor\`eme~13.1; for this ordering of the
factors see C.~Daw and M.~Orr, \emph{Quantitative reduction theory
and unlikely intersections}, \S3.3, Theorem~3.2.

Every period has the form $Z=u\cdot iI$: at each place one can first
apply the Levi factor $Y^{1/2}$ and then translate by the Hermitian
real part $X$. Write $u=\gamma b\omega ak$ using the covering.
Left multiplication by $g=b^{-1}\gamma^{-1}$ sends $Z$ to
$\omega a\cdot iI$. A common $c$ clears the entries of all $b^{-1}$;
the integral lattice stabilizer has integral inverses, so $cg$ is
integral. The inequalities on $A_t$ imply
$a_i\ge t^{h-i}\sqrt t$. At a place the upper-block parabolic matrix
$\omega$ acts by $W\mapsto V_\tau W V_\tau^*+B_\tau$ with
$B_\tau$ Hermitian. The least singular value of $V_\tau$ is bounded
below throughout $\Omega$ and the finite set of places. Hence
$\operatorname{Im}(gZ)_\tau=V_\tau\operatorname{diag}(a_i^2)V_\tau^*$
has a uniform strictly positive lower eigenvalue as claimed.
% END A75B2

% BEGIN A75B3
\paragraph{A75B3. A fixed-level invariant norm for the actual difference.}
Suppose $n$ is even, $4h<n$, and the definite Siegel--Weil identity
holds in degree $h-1$ for every admitted family. B75S4 gives the
boundary limit in each nonzero rational direction in degree $h$
under this induction hypothesis. For fixed positive integers $M,R$ there exists
$C$ such that every $(M,R)$-admitted family $\phi$ bounded in modulus
by $B$ on all level-$M$ lattices satisfies
\[
  |D_\phi(Z)|D(\operatorname{Im}Z)^{n/2}\le CB
  \qquad\text{at every allowed period }Z.
\]
Here $D_\phi$ is the native stabilizer-weighted class sum of the
Hermitian theta series minus the Eisenstein series, both evaluated
at zero elliptic variables.

First work in any fixed high cone $Y\ge\delta I$. Each of the finitely
many chosen lattice representatives $\Lambda$ lies between
$s_\Lambda\mathcal O_E^n$ and $s_\Lambda^{-1}\mathcal O_E^n$ for
some positive integer $s_\Lambda$. Take $R_0=R\prod_\Lambda s_\Lambda$.
The support and period clauses of admission then give level $R_0$
in ordinary integral matrix coordinates for every representative,
uniformly in $\phi$. C75E4 at $\delta/4$ bounds $|D_\phi|$ there by
$2C_0 B\sum_{[\Lambda]}\operatorname{wt}(\Lambda)$. Holomorphy
follows termwise as in O74E5 for these finitely many tests.

The zero-variable difference satisfies
$D_{\omega(g)\phi}(gZ)=j(g,Z)D_\phi(Z)$, where
$j(g,Z)=\prod_\tau\det(C_\tau Z_\tau+D_\tau)^n$ for lower blocks
$(C,D)$ of $g$. For the theta side this is the defining finite Weil
transformation. For the Eisenstein side write the terms for
$\omega(g)\phi$ using $\sigma g^{-1}$ in place of $\sigma$; right
multiplication by $g^{-1}$ permutes the bottom spaces. Weil composition
in the numerator and the automorphy cocycle in the denominator give
the displayed identity, using O74E3 for representative independence
and C75E4 for reindexing. For a translation by $2R_0^2\beta$ with
$\beta$ integral Hermitian, its action on a representative test is
multiplication by a quadratic character. This equals one on its
support: with $u=R_0x$ integral, the phase up to sign is $2\pi i$
times $\sum_\tau\tau(2\operatorname{tr}(u\beta u^*))$, an integer
multiple of $2\pi i$ since the trace entry belongs to
$\mathcal O_{E^+}$. The determinant automorphy factor is one. Thus
$D_\phi$ has period lattice $2R_0^2\operatorname{Herm}_h(\mathcal O_E)$.
Use the rational-direction limits and apply A75B1 with $b=n/2$.
This proves the bound in the high cone with a single constant for all
$(M,R)$-admitted $\phi$.

Take now the uniform $c,\delta$ of A75B2. For a given $Z$ use the
matrix $g$ there. The bounded-denominator transport W75A1--W75A2
makes $\omega(g)\phi$ an $(M,R')$-admitted family of supremum norm
at most $C_1B$, with $R',C_1$ independent of $Z,g,\phi$. Apply the
high-cone bound at this one fixed level $R'$ to $gZ$. If
$Q_\tau=C_\tau Z_\tau+D_\tau$, the unitary block identity gives
$Q_\tau^*\operatorname{Im}(gZ)_\tau Q_\tau=Y_\tau$. Hence
\[
  D(\operatorname{Im}(gZ))^{n/2}|j(g,Z)|=D(Y)^{n/2}.
\]
Substitution of the transformation formula proves the bound in the
whole domain. When either $M$ or $R$ is zero, admission has no
instances because its definition requires both to be positive.
% END A75B3
